\chapter{Endliche Horizonte und Fakultätsschalen}
\label{ch:horizonte}

% Register: D1, D2, S1. Erst Beispiele, dann Definition.

\section{Ein Nachmittag in einer Welt mit 24 Zahlen}

Bevor wir definieren, rechnen wir. Wir nehmen drei Stellen. Die erste
darf 0 oder 1 sein, die zweite 0, 1 oder 2, die dritte 0 bis 3. Die
Gewichte laufen von rechts: Die letzte Stelle zählt Einer, die
mittlere zählt Vierer (denn die letzte Stelle wird mit ihren vier
möglichen Werten „voll“, bevor die mittlere weiterschaltet), die
erste zählt Zwölfer. Wer alle erlaubten Kombinationen der Reihe nach
aufschreibt, zählt von 0 bis 23:

\begin{center}
\begin{tabular}{cccc}
$000=0$ & $001=1$ & $002=2$ & $003=3$\\
$010=4$ & $011=5$ & $012=6$ & $013=7$\\
$020=8$ & $021=9$ & $022=10$ & $023=11$\\
$100=12$ & $101=13$ & $102=14$ & $103=15$\\
$110=16$ & $111=17$ & $112=18$ & $113=19$\\
$120=20$ & $121=21$ & $122=22$ & $123=23$\\
\end{tabular}
\end{center}

Keine Lücke, keine Doppelung: genau $2\cdot3\cdot4=24$ Zahlen. Bei
$123$ ist Schluss -- jede Stelle steht auf ihrem Maximum, und links
anbauen geht nicht, weil es unter der Basis zwei nichts gibt. Diese
Welt mit 24 Plätzen nennen wir den Horizont $H_3$. Wer die 24
erreichen will, braucht eine neue Welt: $H_4$ mit
$2\cdot3\cdot4\cdot5=120$ Plätzen, in der alle Gewichte andere sind.

\section{Die Definitionen}

\begin{definition}[Horizont; Register D1]
Für $n\in\Nn$ sei
$\Hn{n}=\{(a_1,\ldots,a_n)\mid 0\le a_i\le i\}$.
Die Stelle $i$ besitzt die Basis $i+1$; damit ist
$|\Hn{n}|=2\cdot3\cdots(n{+}1)=(n+1)!$.
\end{definition}

\begin{definition}[Wert; Register D2]
$V_n(a)=\sum_{i=1}^{n}a_i\,\dfrac{(n+1)!}{(i+1)!}$, gleichwertig
rekursiv $v_0=0$ und $v_i=(i{+}1)v_{i-1}+a_i$ von links nach
rechts, mit $V_n(a)=v_n$.
\end{definition}

\begin{satz}[Register S1]
$V_n\colon\Hn{n}\to\{0,\ldots,(n+1)!-1\}$ ist bijektiv.
\end{satz}

\begin{proof}
Induktion über $n$. Für $n=0$ ist nichts zu zeigen. Die Rekursion
liefert
\[
V_n(a)=(n{+}1)\,V_{n-1}(a_1,\ldots,a_{n-1})+a_n,
\qquad 0\le a_n\le n;
\]
Division mit Rest durch $n+1$ bestimmt also $a_n$
und den Rest der Ziffern eindeutig (Injektivität). Der größte Wert
ist $(n{+}1)(n!-1)+n=(n{+}1)!-1$, und da Definitionsbereich und
Zielbereich beide $(n+1)!$ Elemente haben, folgt die Surjektivität.
\end{proof}

Die Umkehrabbildung $E_n$ -- einen Wert in seine Ziffern zerlegen --
ist die fortgesetzte Division mit Rest durch $n{+}1, n, \ldots, 2$,
also genau der Algorithmus, mit dem man Sekunden in Tage, Stunden,
Minuten zerlegt, nur mit regelmäßig wachsenden statt willkürlich
gemischten Basen.

\section{Fakultätsschalen}

Jede natürliche Zahl $x$ hat ein kleinstes Zuhause: den
\emph{Minimalhorizont} $\mu(x)=\min\{n:x<(n+1)!\}$. Die natürlichen
Zahlen zerfallen dadurch in \emph{Schalen} $S_n=\{x:\mu(x)=n\}$ --
und die sind begrifflich etwas anderes als die Horizonte: Der
Horizont $\Hn{n}$ ist der ganze Raum und beherbergt die Schalen
$S_0$ bis $S_n$; die Schale $S_n$ ist sein \emph{Neuland}, also die
Werte, die erst mit ihm möglich werden.

\begin{center}
\begin{tabular}{lll}
Schale & Bereich & neue Werte\\
\hline
$S_0$ & nur die $0$ & $1$\\
$S_1$ & nur die $1$ & $1$\\
$S_2$ & $2$ bis $5$ & $4$\\
$S_3$ & $6$ bis $23$ & $18$\\
$S_4$ & $24$ bis $119$ & $96$\\
$S_5$ & $120$ bis $719$ & $600$\\
$S_6$ & $720$ bis $5039$ & $4320$\\
\end{tabular}
\end{center}

Die Schalengrenzen sind zahlentheoretisch markante Orte: Für
$n\ge1$ ist die untere Grenze der Schale $S_n$ die Fakultät $n!$ --
durch \emph{jede} Zahl von $1$ bis $n$ teilbar. Die Zahl direkt
davor, $n!-1$, ist durch \emph{keine} Zahl von $2$ bis $n$ teilbar.
Der Schritt von $719$ auf $720$ wechselt also mit einem einzigen
Nachfolger von maximaler Teilerfremdheit zu maximaler Teilbarkeit --
und genau an dieser Stelle erzwingt die Horimetrik den Umzug in eine
neue Welt.

\begin{beispiel}
Die Zahl $7$ wohnt minimal in $H_3$ (denn $6\le7<24$) und schreibt
sich dort $013$. In $H_4$ schreibt sie sich $0012$, in $H_5$ als
$00011$, in $H_6$ als $000010$ -- und ab $H_7$ stabil als
$0\cdots07$: Sobald der Horizont groß genug ist, passt die ganze
Zahl in die letzte Ziffer, und weitere führende Nullen sind endlich
wirklich harmlos. Diese Lebensgeschichte erzählen wir später in
einem eigenen Interludium in Ruhe.
\end{beispiel}

Damit steht das Fundament: endliche Welten mit Fakultätsgröße, eine
Bijektion zwischen Ziffern und Werten, und Schalen mit markanten
Grenzen. Was noch fehlt, ist der eigentliche Hauptdarsteller -- die
zweite Identität jeder Zahl. Sie betritt im nächsten Kapitel die
Bühne.
